\begin{abstract}

% Springer asks for 150--250 words. ~215 here.
%
% Structure:
%   1. Setting      -- four-sequence MRI, sequences often missing in practice
%   2. Gap          -- existing methods keep accuracy but return one mask
%   3. Method       -- 3D latent diffusion, modality dropout, N samples
%   4. Contribution -- the restricted uncertainty average
%   5. Evaluation   -- 2601 studies, accuracy holds, uncertainty responds
%
% !! The Dice figure below is marked XX. As of the current run whole-tumour
%    Dice is about 92 on the test split, but the baseline comparison is still
%    being re-run, so insert the final number deliberately rather than letting
%    a provisional one reach submission.

Brain tumour segmentation from magnetic resonance imaging normally assumes
that all four imaging sequences are available, but in clinical practice one or
more are frequently missing. Existing methods for incomplete inputs maintain
accuracy by learning modality-invariant representations or by distilling from
a full-modality teacher, yet they return a single segmentation mask regardless
of how much information the input contains. A prediction made from one
sequence is therefore presented in the same form as one made from four, and
the reader is given no indication of how confident the model is. This paper
proposes a three-dimensional conditional latent diffusion framework that
produces several plausible segmentations for a given study, together with a
voxel-wise map of their disagreement. Two volumetric autoencoders compress
images and annotations before diffusion, which makes sampling over $128^{3}$
volumes practical, and modality availability is randomly resampled during
training so that a single model serves all fifteen acquisition
configurations. We further show that averaging predictive variance over the
whole volume conflates confidence with lesion size, and propose a restricted
average that remains comparable as sequences are withdrawn. Evaluations on
$2{,}601$ studies from BraTS 2023 and 2024 indicate that the method maintains
whole-tumour accuracy under incomplete inputs, and that the predicted
uncertainty increases as sequences are removed.

\keywords{Brain Tumour Segmentation \and Latent Diffusion Model \and
Missing Modalities \and Uncertainty Estimation \and Magnetic Resonance
Imaging}
\end{abstract}
